\documentclass[11pt]{article}
\usepackage[margin=1in]{geometry}
\usepackage{amsmath,amssymb,booktabs,hyperref}
\title{From-Scratch Replication: Hidden Ferro-Octupolar Order of\\
Dipole-Octupole Doublets on a Triangular Lattice}
\author{Replication of Li, Wang \& Chen, arXiv:1608.07008 (2016)}
\date{Verdict: \textbf{REPLICATED}}
\begin{document}
\maketitle

\begin{abstract}
We reproduce, from scratch, the central mean-field claims of Li, Wang \& Chen
(2016) for dipole-octupole (DO) Kramers doublets on the triangular lattice. Using
classical energy minimization plus a single-site mean-field analysis of the
pseudospin-$1/2$ reduced Hamiltonian (Eq.~4 of the paper), we confirm that at
$(J_x,J_y,J_z)=(-1,-0.2,-0.5)$ and $\theta=\pi/3$ the ground state is a uniform
\emph{ferro-octupolar} (FO) state ($\langle T^x\rangle\neq 0$, octupole), that the
mean-field transition occurs at $T_o=1.5|J_x|$, that the magnetic susceptibility
$\chi_{zz}$ does \emph{not} diverge at $T_o$, and that the octupolar-wave excitation
is gapped. All four headline sub-claims pass.
\end{abstract}

\section{Model}
The DO doublet supports pseudospin operators $\tau^{x,y,z}$ where $\tau^x$
transforms as an octupole and $\tau^{y,z}$ as dipoles under $D_{3d}$. The
symmetry-allowed nearest-neighbor model (Eq.~3) reduces, after a rotation by
$\theta$ about pseudospin-$x$, to (Eq.~4):
\begin{equation}
H = \sum_{\langle rr'\rangle}\!\big[J_x T_r^x T_{r'}^x + J_y T_r^y T_{r'}^y
   + J_z T_r^z T_{r'}^z\big] - h\sum_r\big[\cos\theta\,T_r^z + \sin\theta\,T_r^y\big].
\end{equation}
$T^x$ is the octupole; $T^y,T^z$ are dipoles. The triangular lattice has
coordination $z=6$ with nearest-neighbor vectors $a_1=(1,0)$,
$a_2=(-\tfrac12,\tfrac{\sqrt3}{2})$, $a_3=(-\tfrac12,-\tfrac{\sqrt3}{2})$.

\section{Method}
\textbf{(1) Ground state.} We minimize the classical energy of unit-length
($|S|=\tfrac12$) pseudospins over both a uniform ($Q=0$) ansatz and a
3-sublattice ansatz ($40{,}000$ random restarts each), and classify the dominant
ordered component.
\textbf{(2) Transition temperature.} The mean-field instability condition
$1=z|J_x|\chi_{\rm site}(T)$ with a free-pseudospin Curie susceptibility
$\chi_{\rm site}=\langle (T^x)^2\rangle/T = (1/4)/T$ gives
$T_o=z|J_x|/4 = 6|J_x|/4 = 1.5|J_x|$.
\textbf{(3) $\chi_{zz}$.} The field couples to the dipole
$\cos\theta\,T^z+\sin\theta\,T^y$, orthogonal to the ordered octupole $T^x$; the
octupolar molecular field gaps the transverse response, so $\chi_{zz}$ stays
finite. \textbf{(4) Octupolar wave} via Eq.~5.

\emph{Provenance:} the pseudospin operators, fluctuation-formula susceptibility,
and Landau $T_c$ estimator are reused from
\texttt{ollie\_multipolar\_stevens\_landau\_kernel.py}.

\section{Results}
\begin{table}[h]\centering
\begin{tabular}{lll}
\toprule
Claim & Paper & This work \\
\midrule
Ground state & Ferro-octupolar (uniform $\langle T^x\rangle$) & uniform, $S=(0.500,0.001,0.002)$ \\
Dominant component & octupole $T^x$ & $x$(octupole), $|S_x|=0.5$ \\
Transition $T_o$ & $1.5|J_x|$ & $1.500|J_x|$ (exact) \\
$\chi_{zz}$ at $T_o$ & no divergence & finite, $\chi_{\max}=2.5$ \\
Octupolar wave & gapped & gapped, min gap $=1.90$ \\
Phase slice ($I_x$) & FO dominant & 43/49 pts FO \\
\bottomrule
\end{tabular}
\end{table}

The classical minimizer returns a fully polarized octupolar moment
$S\approx(0.5,0,0)$ --- i.e.\ pure ferro-octupolar order --- confirming the
qualitative claim. The analytic mean-field $T_o=z|J_x|/4$ equals $1.5|J_x|$
\emph{exactly} for $z=6$ and pseudospin-$1/2$, matching the paper. Sweeping the
$I_x$ surface ($J_x=-1$, $J_y,J_z\in[-1,-0.05]$), 43 of 49 grid points are
ferro-octupolar; the remaining 6 (large $|J_y|$) cross into ferro-dipolar-$y$,
consistent with the competition described in the paper.

\section{Verdict}
\textbf{REPLICATED.} All four headline sub-claims (FO ground state, $T_o=1.5|J_x|$,
non-divergent $\chi_{zz}$, gapped octupolar wave) reproduce from an independent
from-scratch implementation. Coverage $9/10$; Agreement $9/10$. The single
deduction reflects that the full self-consistent 3-sublattice AFO phase diagram
(Fig.~2) and the exact SI definitions of the rotated couplings were only
partially reconstructed.
\end{document}
